% !TeX root = ../navier-stokes-manuscript.tex
% ============================================================
% 2.4 The Spatial Column Inside Critical Height
% ============================================================

\subsection{The Spatial Column Inside Critical Height}\label{sub:TheSpatialColumnInsideCriticalHeight}

For \(s\geq0\), set
\[
  E_s(t)
  :=
  \frac12\norm{\Lambda^su(t)}_2^2.
\]
The critical height is \(H=E_{1/2}\).
To see what one time derivative of \(H\) contains, keep \(s\) arbitrary and derive the energy balance at height \(s\).

Apply \(\Lambda^s\) to the momentum equation and take the \(L^2\) inner product with \(\Lambda^su\).
The time derivative gives
\[
  \pair{\Lambda^su}{\Lambda^s\partial_tu}
  =
  \frac{d}{dt}
  \frac12\norm{\Lambda^su}_2^2
  =
  E_s'.
\]
Rapid decay and the commutation of \(\Lambda^s\) with spatial derivatives give
\[
  \nu\pair{\Lambda^su}{\Lambda^s\Delta u}
  =
  -\nu\norm{\Lambda^{s+1}u}_2^2
  =
  -2\nu E_{s+1}.
\]
The pressure pairing vanishes in this global balance:
\[
  \pair{\Lambda^su}{\Lambda^s\nabla p}
  =
  -\pair{\nabla\cdot\Lambda^su}{\Lambda^sp}
  =
  0.
\]
For the remaining signed transfer, write
\[
  \mathcal P_s(t)
  :=
  -\pair{\Lambda^su}
  {\Lambda^s\bigl((u\cdot\nabla)u+\nabla p\bigr)}.
\]
The pressure term is retained in this definition because \(\mathcal P_s\) comes from the complete momentum equation, even though its global pairing is zero.
The exact balance is
\[
  E_s'
  =
  \mathcal P_s-2\nu E_{s+1}.
\]
One time derivative of the spatial energy at height \(s\) contains the energy one full Sobolev step higher.
It also contains the signed transfer \(\mathcal P_s\) generated by the lower rows of the mixed recurrence.

Put \(s=1/2\).
The first critical-height row is
\[
  H'
  =
  \mathcal P_{1/2}
  -2\nu E_{3/2}.
\]
The next row requires the balance at \(s=3/2\),
\[
  E_{3/2}'
  =
  \mathcal P_{3/2}
  -2\nu E_{5/2}.
\]
Substituting it into the derivative of \(H'\) gives
\begin{align*}
  H''
  &=
  \mathcal P_{1/2}'
  -2\nu E_{3/2}'
  \\
  &=
  \mathcal P_{1/2}'
  -2\nu\mathcal P_{3/2}
  +(2\nu)^2E_{5/2}.
\end{align*}
The third derivative uses
\[
  E_{5/2}'
  =
  \mathcal P_{5/2}
  -2\nu E_{7/2}
\]
and gives
\[
  H'''
  =
  \mathcal P_{1/2}''
  -2\nu\mathcal P_{3/2}'
  +(2\nu)^2\mathcal P_{5/2}
  -(2\nu)^3E_{7/2}.
\]
One more differentiation gives
\[
  H^{(4)}
  =
  \mathcal P_{1/2}'''
  -2\nu\mathcal P_{3/2}''
  +(2\nu)^2\mathcal P_{5/2}'
  -(2\nu)^3\mathcal P_{7/2}
  +(2\nu)^4E_{9/2}.
\]
The highest spatial energy advances by one Sobolev step each time, while every transfer encountered on the way remains in the row.

The pattern follows by induction.
For every integer \(n\geq1\),
\[
  H^{(n)}
  =
  (-2\nu)^nE_{n+1/2}
  +
  \sum_{j=0}^{n-1}
  (-2\nu)^j
  \partial_t^{\,n-1-j}\mathcal P_{j+1/2}.
\]
Indeed, differentiating the \(n\)-th row differentiates each existing transfer term.
Replacing
\[
  E_{n+1/2}'
  =
  \mathcal P_{n+1/2}
  -2\nu E_{n+3/2}
\]
creates the next transfer term and the next highest spatial energy with exactly the displayed coefficient.

The first rows keep their distinct quantities intact:
\[
  \begin{array}{c|c}
    \text{temporal row}
    &
    \text{highest spatial energy appearing in that complete row}
    \\
    \hline
    H
    &
    E_{1/2}
    \\
    H'
    &
    -2\nu E_{3/2}
    \\
    H''
    &
    (2\nu)^2E_{5/2}
    \\
    H'''
    &
    -(2\nu)^3E_{7/2}
    \\
    H^{(4)}
    &
    (2\nu)^4E_{9/2}
  \end{array}
\]
For example, \(H'\) is the full quantity \(\mathcal P_{1/2}-2\nu E_{3/2}\).
The table identifies the highest spatial term inside each complete temporal row.

The energies \(E_s\) use the homogeneous Sobolev scale because that scale makes the dilation exact.
Continuation is usually stated in the inhomogeneous scale.
For every \(s>0\),
\[
  \norm{u(t)}_{H^s}^2
  \simeq
  \norm{u(t)}_2^2
  +
  \norm{\Lambda^su(t)}_2^2
  =
  \norm{u(t)}_2^2+2E_s(t).
\]
The unforced energy identity gives
\[
  \frac12\norm{u(t)}_2^2
  +
  \nu\int_0^t\norm{\nabla u(\tau)}_2^2\,d\tau
  =
  \frac12\norm{u_0}_2^2,
\]
so the low-frequency term is already controlled by the datum.
A bound for \(E_s\) supplies the corresponding inhomogeneous \(H^s\) bound, and an \(L^1\)-in-time bound for \(E_s\) supplies the high-frequency part of an \(L^2_tH^s_x\) bound on every finite interval.

Finite remaining time pushes the threat upward through the temporal tower: once one rung escapes, every rung above it must escape too.
The critical-height identities expose the opposite spatial movement: each derivative of \(H\) contains a higher Sobolev energy, while a singular endpoint would have to lose spatial regularity.
These two directions are balanced and opposing.
Any finite number of temporal rows reaches only one finite spatial height, so keeping both directions on the same velocity field requires the intersection of all finite rows.
